\begin{otherlanguage}{english}
\section*{Executive summary}
\addcontentsline{toc}{section}{Executive summary}

\textbf{Context.} This work was carried out during the end-of-studies internship at
\acompleter{company name}, in \acompleter{team}, from \acompleter{dates}. It deals with the valuation of options on
French government bonds (OATs). A call option gives the right to buy the bond at a future date for a price fixed
today; a put option gives the right to sell it. The contract studied is the one traded by banks: the option is
collateralised, with collateral paid at the €STR overnight rate, and its forward price follows from the cost of
financing the bond in the repurchase (repo) market.

\textbf{Research question.} The price of an OAT depends on two sources of risk: the level of risk-free rates, and
the spread between the OAT yield and those rates, which compensates for the credit risk of the French State, the
liquidity of the bond and supply effects. Standard interest rate models only capture the first source, and there
is no liquid market for OAT options from which the volatility of the spread could be inferred. The question is
therefore twofold: should the spread be modelled as a separate risk factor to price an OAT option, and can its
parameters be identified from the available data?

\textbf{Approach.} The model, taken from the literature (Brigo and Mercurio, 2006; Russo, Giacometti and Fabozzi),
describes the OAT short rate as the sum of two Gaussian factors: the risk-free rate and the spread. The rate factor
is calibrated to swaption prices, which are very liquid interest rate options. For the spread factor, we first
show that the calibration procedure proposed in the reference paper, based on the OAT yield curve alone, cannot
identify any parameter. The spread parameters are therefore estimated on six years of daily data (2021-2026),
with two econometric methods taught at ENSAE: the generalised method of moments and maximum likelihood with the
Kalman filter. The option price is computed exactly, and the computation is validated by Monte Carlo simulation.
Market data come from Bloomberg (8 September 2026); the option is written on the 4.10\% OAT maturing in 2046, with
expiries from one to ten years.

\textbf{Results.} Four results stand out.
\begin{itemize}
\item The volatility of the spread is about 34 basis points per year and its correlation with rates is 0.13; it
switches between a calm regime (23 basis points) and turbulent episodes (51 basis points), such as 2022 or the
French political crises of 2024, but these episodes are too short to affect the price.
\item Accounting for the spread makes the option 17\% more expensive, at all expiries.
\item For an option on a single OAT, this effect amounts to a higher volatility: a one-factor model whose
volatility is raised by about 17\% reproduces the prices within 0.05\%.
\item The closed-form approximation of the reference paper is accurate for at-the-money options but errs by 3
to 4\% for options far from the money: the exact price should be preferred.
\end{itemize}
Finally, for long expiries, the repo financing cost weighs as much on the price as the spread itself.

\textbf{Limitations and extensions.} For lack of quotes, the repo curve used is fictitious; a market curve would
replace it without changing the computations. Statistical tests show that a single factor per curve does not
capture all the deformations of the spread, and that the correlation between spread and rates is unstable over
time. These findings call for richer models, mostly relevant for products sensitive to the spread between
several bonds.
\end{otherlanguage}
